\documentclass[11pt,a4paper]{article}

\usepackage[utf8]{inputenc}
\usepackage[T1]{fontenc}
\usepackage[french]{babel}
\usepackage{amsmath,amssymb,amsthm,mathtools}
\usepackage{geometry}
\usepackage{enumitem}
\usepackage{booktabs}
\usepackage{array}
\usepackage{longtable}
\usepackage{hyperref}
\usepackage{microtype}
\usepackage{stmaryrd}
\usepackage{fancyhdr}
\usepackage{parskip}

\geometry{margin=2.5cm}
\setlist[itemize]{leftmargin=1.7em}
\setlist[enumerate]{leftmargin=1.9em}
\setlength{\headheight}{14pt}
\hypersetup{colorlinks=true,linkcolor=black,urlcolor=blue,citecolor=black}
\pagestyle{fancy}
\fancyhf{}
\lhead{Infini vs fini -- Version 2}
\rhead{Racine et inracine}
\cfoot{\thepage}

\title{\textbf{Infini vs fini -- Formalisation rigoureuse d'une hypothèse}\\
\large Version 2 : racines, inracines, contraintes et récursivité}
\author{}
\date{}

\newtheorem{definition}{Définition}[section]
\newtheorem{axiom}{Axiome}[section]
\newtheorem{proposition}{Proposition}[section]
\newtheorem{lemma}{Lemme}[section]
\newtheorem{remark}{Remarque}[section]
\newtheorem{consequence}{Conséquence}[section]
\newtheorem{question}{Question ouverte}[section]
\newtheorem{warning}{Point de vigilance}[section]
\newtheorem{example}{Exemple}[section]

\newcommand{\Rac}{\mathsf{R}}
\newcommand{\Inrac}{\mathsf{I}}
\newcommand{\Cadre}{\mathsf{K}}
\newcommand{\Con}{\mathsf{C}}
\newcommand{\Sol}{\mathcal{S}}
\newcommand{\Det}{\mathsf{D}}
\newcommand{\Actual}{\mathsf{A}}
\newcommand{\Per}{\mathsf{P}}
\newcommand{\Ex}{\mathsf{Ex}}
\newcommand{\Refi}{\mathsf{Ref}}
\newcommand{\ClotureOnt}{\operatorname{ClotureOntologique}}
\newcommand{\Loc}{\mathsf{Loc}}
\newcommand{\sat}{\models}
\newcommand{\poss}{\Diamond}
\newcommand{\nextlev}{\leadsto}

\begin{document}
\maketitle

\begin{abstract}
Ce document propose une révision complète de l'hypothèse développée dans la
Version 1. La notion d'\emph{inracine} devient l'outil conceptuel central pour
remplacer la formulation directe selon laquelle « l'infini contraint le fini ».
L'objectif est de préserver l'intuition originale --- un fond non clôturable,
des déterminations locales, des contraintes, des solutions et une récursivité
entre niveaux --- tout en évitant deux confusions : assimiler l'infini à un
objet gigantesque, et assimiler l'inracine à un simple contenu encore inconnu.

La théorie distingue ainsi :
\[
\boxed{
\text{inracine}
\;\leftrightarrow\;
\text{dépassement d'un cadre}
\;\leftrightarrow\;
\text{racine locale}
}
\]
et propose une structure récursive générale :
\[
\boxed{
\Inrac_0\;
extarrow\;\Cadre_0\;\to\;\Rac_0\;\to\;
\Inrac_1\;\to\;\Cadre_1\;\to\;\Rac_1\;\to\;\cdots
}
\]
Le symbole de flèche n'est pas interprété par défaut comme une causalité
physique. Il exprime selon le contexte une relation de conditionnement,
de détermination, de dérivation ou de passage de niveau.

La Version 2 est une formalisation métaphysique et méta-théorique. Elle ne
démontre pas que le réel physique est effectivement infini, qu'il existe une
infinité actuelle d'objets, ni que l'univers est un ordinateur. Elle fournit en
revanche un langage rigoureux pour exprimer une thèse ontologique : une
\emph{racine} est une détermination locale obtenue sous certaines contraintes,
et une \emph{inracine} désigne le dépassement non exhaustivement clôturable du
cadre dans lequel cette racine est déterminée.
\end{abstract}

\section{Principe de fidélité et statut de la théorie}

\begin{definition}[Principe de fidélité]
Une théorie doit d'abord être reconstruite selon ses propres primitives avant
d'être comparée aux définitions standards de la métaphysique, des
mathématiques, de la physique ou de la théorie de la connaissance.
\end{definition}

La Version 2 poursuit exactement trois opérations distinctes :
\begin{enumerate}
    \item expliciter les primitives de l'hypothèse ;
    \item dériver les conséquences internes à ces primitives ;
    \item comparer ensuite ces conséquences à des cadres externes.
\end{enumerate}

\begin{warning}[Pas de tribunal ontologique préalable]
Une définition mathématique interne à un système formel établit ce qui vaut
dans ce système. Elle ne constitue pas, à elle seule, une démonstration que
la structure ainsi définie est une catégorie ontologique ultime du réel.
Réciproquement, contester une interprétation ontologique ne revient pas à
contester la cohérence des mathématiques qui utilisent cette définition.
\end{warning}

La théorie proposée est donc de nature \emph{ontologique et relationnelle}.
Elle ne cherche pas d'abord à supprimer le langage du fini, mais à lui retirer
le statut de primitive ontologique absolue.

\section{Du problème « infini vs fini » au couple « racine vs inracine »}

La Version 1 partait de l'idée :
\[
\text{fond non limité}
\to
\text{contraintes}
\to
\text{problèmes}
\to
\text{solutions}.
\]
La Version 2 introduit une distinction supplémentaire : toute solution ou
détermination locale peut être considérée comme une \emph{racine}, relativement
à un cadre, tandis que ce qui dépasse la clôture de ce cadre constitue son
\emph{inracine}.

On passe donc de la simple opposition :
\[
\text{infini}\;\text{vs}\;\text{fini}
\]
à une architecture asymétrique :
\[
\boxed{
\text{racine}\;\text{vs}\;\text{inracine}
}
\]
où la relation fondamentale n'est pas celle de deux objets symétriques, mais
celle d'une détermination locale et de son dépassement de cadre.

\subsection{Pourquoi le terme « inracine » ?}

Le mot \emph{inracine} est ici introduit comme un terme théorique original. Il
est construit comme l'envers conceptuel de « racine », non comme un terme déjà
standardisé d'une discipline particulière.

La racine évoque ici ce qui est suffisamment déterminé pour servir de point
d'appui local. L'inracine désigne ce qui, relativement à ce point d'appui,
ne peut être réduit à sa clôture locale.

\begin{remark}
L'inracine ne signifie donc pas simplement « extérieur », « plus grand »,
« caché », « inconnu », « imaginaire » ou « infini ». La notion est plus précise :
elle désigne une \emph{structure de non-clôture relative à un cadre}.
\end{remark}

\section{Primitives formelles}

Afin d'éviter de transformer l'ontologie en une simple théorie des ensembles,
les primitives sont typées.

\begin{definition}[Inracine]
Pour un cadre local \(\Cadre\), son \emph{inracine} \(\Inrac(\Cadre)\) est la
structure théorique représentant ce qui excède la capacité de clôture du cadre
\(\Cadre\), sans que cet excédent soit lui-même identifié à un objet extérieur
ou à une totalité exhaustive.
\end{definition}

Il est utile d'écrire :
\[
\boxed{
\Inrac(\Cadre)=\text{dépassement non clôturé du cadre }\Cadre.
}
\]
Le terme « dépassement » est relationnel : il est toujours le dépassement
d'un cadre déterminé.

\begin{definition}[Cadre]
Un \emph{cadre} \(\Cadre\) est une structure de détermination qui spécifie au
moins :
\begin{enumerate}
    \item un domaine local de candidats ou d'états ;
    \item un langage ou un mode de représentation ;
    \item un ensemble de contraintes \(\Con\) ;
    \item éventuellement des ressources, échelles et règles de sélection.
\end{enumerate}
\end{definition}

\begin{definition}[Contrainte]
Une \emph{contrainte} est une condition qui restreint, structure ou organise
l'espace des déterminations admissibles dans un cadre donné.
\end{definition}

La contrainte n'est donc pas nécessairement une cause. Elle peut être physique,
causale, logique, géométrique, matérielle, cognitive, computationnelle,
linguistique, mathématique, sociale, informationnelle ou conceptuelle.

\begin{definition}[Espace local de solutions]
Pour un cadre \(\Cadre\) possédant un espace de candidats \(\mathcal D_{\Cadre}\)
et un ensemble de contraintes \(\Con_{\Cadre}\), on définit :
\[
\Sol(\Cadre)=
\{D\in\mathcal D_{\Cadre}\mid D\sat\Con_{\Cadre}\}.
\]
Une détermination admissible est donc un élément de l'espace de solutions.
\end{definition}

\begin{definition}[Racine]
Une \emph{racine} \(\Rac\) est une détermination locale distinguée relativement
à un cadre \(\Cadre\). Elle peut être :
\begin{enumerate}
    \item une solution effectivement sélectionnée ;
    \item une solution construite ;
    \item une valeur observée ;
    \item une structure stabilisée ;
    \item ou une représentation suffisamment déterminée pour servir de point
de départ à un niveau ultérieur.
\end{enumerate}
\end{definition}

La notation :
\[
\boxed{\Rac\in\Sol(\Cadre)}
\]
exprime le cas où la racine est effectivement compatible avec les contraintes.

\begin{definition}[Localité]
Une détermination est dite \emph{locale} lorsqu'elle est déterminée relativement
à un cadre donné. La localité ne désigne donc pas nécessairement une petite
région spatiale ; elle peut être spatiale, temporelle, causale, structurelle,
informationnelle, perceptive, mathématique ou conceptuelle.
\end{definition}

Ainsi :
\[
\boxed{
\text{local}=\text{déterminé relativement à un cadre donné}.
}
\]

\section{Axiomes de la Version 2}

\begin{axiom}[Fond non clôturable]
Il existe un principe fondamental \(\mathfrak I_\ast\), appelé ici \emph{inracine
fondamentale}, tel qu'il n'est pas lui-même une racine résultant d'une clôture
ontologique supérieure.
\end{axiom}

Cet axiome remplit le rôle ontologique qui était auparavant attribué directement
à « l'infini ». Il est important de ne pas le confondre avec un objet maximal.

\begin{axiom}[Non-totalisation]
Il n'existe pas, dans la théorie, de cadre local \(\Cadre^{\ast}\) qui puisse être
identifié à une clôture exhaustive de \(\mathfrak I_\ast\).
\end{axiom}

On peut schématiser :
\[
\boxed{
\neg\exists \Cadre^{\ast}\;\text{tel que}\;
\Cadre^{\ast}\text{ clôture exhaustivement }\mathfrak I_\ast.
}
\]
Ce n'est pas un énoncé sur la cardinalité d'un ensemble ; c'est une condition
méta-ontologique de non-clôture.

\begin{axiom}[Détermination sous contraintes]
Toute racine locale pertinente pour la théorie est relative à un cadre de
contraintes :
\[
\Rac\in\Sol(\Cadre).
\]
La théorie ne pose donc pas la racine comme une entité ontologiquement
indépendante du cadre de détermination.
\end{axiom}

\begin{axiom}[Récursivité]
Lorsqu'une racine devient suffisamment déterminée pour servir de nouvelle base
de détermination, elle peut contribuer à la formation d'un cadre ultérieur.
On écrit :
\[
\Rac_n\nextlev\Cadre_{n+1}.
\]
Cette possibilité n'implique pas qu'une racine engendre nécessairement une
seule et unique contrainte ultérieure.
\end{axiom}

\begin{axiom}[Relativité de la clôture]
Une clôture locale est toujours relative à son cadre. Le passage de \(\Cadre_n\)
à \(\Cadre_{n+1}\) peut révéler des déterminations qui n'étaient pas pertinentes,
accessibles ou exprimables dans \(\Cadre_n\).
\end{axiom}

\section{L'inracine n'est pas l'inconnu}

Cette distinction est indispensable à la cohérence du concept.

\begin{definition}[Inracine épistémique]
L'inracine épistémique d'un cadre \(\Cadre\), notée \(\Inrac^{\mathrm{epi}}(\Cadre)\),
désigne ce qui dépasse les capacités actuelles de connaissance,
d'observation, de description ou de calcul de l'agent considéré.
\end{definition}

\begin{definition}[Inracine ontologique]
L'inracine ontologique \(\Inrac^{\mathrm{ont}}(\Cadre)\) désigne ce qui excède
la clôture ontologique du cadre \(\Cadre\), indépendamment du fait que ce
surplus soit actuellement connu ou inconnu.
\end{definition}

\begin{proposition}[Non-équivalence]
En général :
\[
\boxed{
\Inrac^{\mathrm{epi}}(\Cadre)\neq
\Inrac^{\mathrm{ont}}(\Cadre).
}
\]
\end{proposition}

\begin{proof}
Un domaine peut être parfaitement connu par un agent tout en restant local et
non exhaustif ontologiquement. À l'inverse, un agent peut échouer à connaître
un contenu qui appartient à un cadre fermé par stipulation. L'échec cognitif
ne suffit donc pas à définir la non-clôture ontologique.
\end{proof}

La théorie peut utiliser l'inracine épistémique pour parler de perception,
mais son noyau ontologique doit reposer sur l'inracine ontologique.

\section{Le principe « toute racine porte son inracine »}

Le principe central de la Version 2 est :
\[
\boxed{
\text{Toute racine locale porte son inracine relativement au cadre qui la détermine.}
}
\]

Cela signifie que lorsqu'une racine \(\Rac_n\) est déterminée dans \(\Cadre_n\),
on ne doit pas automatiquement identifier \(\Rac_n\) à une clôture ontologique
ultime.

On introduit une relation binaire :
\[
\Ex(\Rac_n,\Cadre_n)
\]
qui signifie « le cadre \(\Cadre_n\), relativement à la racine \(\Rac_n\),
présente un dépassement non clôturé ».

La définition de l'inracine peut alors s'écrire sous forme relationnelle :
\[
\boxed{
\Inrac_n=\Inrac(\Rac_n,\Cadre_n)
\quad\text{si}\quad
\Ex(\Rac_n,\Cadre_n).
}
\]

\begin{remark}
Cette notation ne signifie pas que \(\Inrac_n\) est nécessairement un objet,
un ensemble ou une région ontologique située « derrière » \(\Rac_n\). Elle
représente d'abord une structure relationnelle de non-clôture.
\end{remark}

\section{La structure récursive complète}

La structure de base devient :
\[
\boxed{
\mathfrak I_0
\rightsquigarrow
\Cadre_0
\to
\Sol(\Cadre_0)
\to
\Rac_0
\rightsquigarrow
\Inrac_0
\rightsquigarrow
\Cadre_1
\to
\Sol(\Cadre_1)
\to
\Rac_1
\rightsquigarrow
\Inrac_1
\rightsquigarrow\cdots
}
\]

Une écriture plus compacte est :
\[
\boxed{
(\Cadre_n,\Rac_n,\Inrac_n)
\nextlev
(\Cadre_{n+1},\Rac_{n+1},\Inrac_{n+1}).
}
\]

On peut décomposer chaque transition :
\[
\Cadre_n
\overset{\Con_n}{\longrightarrow}
\Sol(\Cadre_n)
\overset{\Actual_n}{\longrightarrow}
\Rac_n
\overset{\Ex_n}{\longrightarrow}
\Inrac_n
\overset{\Refi_n}{\longrightarrow}
\Cadre_{n+1}.
\]

Ici :
\begin{itemize}
    \item \(\Con_n\) désigne la structure de contraintes ;
    \item \(\Actual_n\) représente, lorsque nécessaire, la sélection ou
    l'actualisation d'une solution ;
    \item \(\Ex_n\) produit l'identification du dépassement relatif ;
    \item \(\Refi_n\) représente une éventuelle reformulation ou extension du
    cadre à partir de ce dépassement.
\end{itemize}

\begin{warning}[Pas d'identité obligatoire entre niveaux]
La théorie ne dit pas que \(\Rac_n=\Rac_{n+1}\). Elle dit qu'une racine peut
servir de point d'appui pour une nouvelle détermination. Les niveaux sont
reliés, mais ne sont pas identiques.
\end{warning}

\section{Racine comme solution locale}

La Version 1 proposait :
\[
\text{solution}=\text{détermination compatible avec des contraintes}.
\]
La Version 2 précise cette idée :
\[
\boxed{
\Rac=\text{solution locale relativement à un cadre de contraintes}.
}
\]

Il ne faut toutefois pas identifier toute racine à une solution mathématique.
Une racine peut être un état physique stabilisé, une perception, une structure,
un résultat expérimental, une configuration causale ou une valeur numérique.

La relation fondamentale reste :
\[
\Rac\in\Sol(\Cadre).
\]

\subsection{Le fini comme prédicat relationnel secondaire}

On remplace la proposition primitive :
\[
D\in\mathrm{Fini}
\]
par une qualification relative :
\[
\boxed{\mathrm{Fini}_{\Cadre}(D)}.
\]

Cette notation signifie : « \(D\) se présente comme borné ou déterminé dans le
cadre \(\Cadre\) ».

Elle ne signifie pas : « \(D\) est absolument borné dans toute ontologie
possible ».

On peut alors poser :
\[
\boxed{
\mathrm{Fini}_{\Cadre}(\Rac)
\quad\not\Rightarrow\quad
\ClotureOnt(\Rac).
}
\]

\begin{consequence}[Déplacement de la charge ontologique]
La question « pourquoi existe-t-il des objets finis ? » est remplacée, au niveau
fondamental, par « quelles structures de contraintes rendent possible une
racine localement bornée, et comment son inracine est-elle définie ? ».
\end{consequence}

\section{Le rôle de la perception}

La perception fournit un cas privilégié de localité, mais elle n'est pas la
source de l'ontologie.

On distingue :
\[
\underbrace{
\text{réalité}
\to
\Cadre_{\mathrm{phys}}
\to
\Rac_{\mathrm{phys}}
}_{\text{niveau ontologique / physique}}
\]
et :
\[
\underbrace{
\Rac_{\mathrm{phys}}
\to
\Cadre_{\mathrm{cogn}}
\to
\Rac_{\mathrm{cogn}}.
}_{\text{niveau épistémique}}
\]

La perception peut alors être dite « locale » parce qu'elle sélectionne une
détermination selon des capacités sensorielles, attentionnelles,
linguistiques et conceptuelles données.

\begin{definition}[Perception locale]
Une perception est locale lorsqu'elle constitue une détermination relative à
un ensemble de capacités et de contraintes de représentation déterminées.
\end{definition}

Cette définition autorise plusieurs couches :
\[
\mathfrak I
\to
\Rac_{\mathrm{phys}}
\to
\Inrac_{\mathrm{phys}}
\to
\Rac_{\mathrm{cogn}}
\to
\Inrac_{\mathrm{cogn}}
\to\cdots
\]
mais les flèches ne doivent pas être lues comme une chaîne de perception
humaine uniquement. La même logique peut être appliquée à des systèmes
physiques, biologiques, computationnels ou mathématiques.

\section{Conséquences ontologiques}

\subsection{1. Non-fondamentalité du fini}

\begin{proposition}
Dans la théorie, le caractère localement fini d'une racine ne suffit pas à lui
conférer un statut de primitive ontologique.
\end{proposition}

\begin{proof}
Par définition, une racine est relative à un cadre et satisfait ses contraintes.
Sa finitude locale est donc un prédicat relatif à ce cadre. Un prédicat relatif
ne peut être utilisé comme primitive absolue sans ajouter un axiome indépendant
établissant cette absoluité.
\end{proof}

\subsection{2. Dépendance relationnelle de la racine}

Si :
\[
\Rac\in\Sol(\Cadre),
\]
alors l'identité théorique de \(\Rac\) dépend au moins partiellement de la
structure de \(\Cadre\).

Cela ne signifie pas que la racine est « seulement subjective ». Cela signifie
que sa détermination locale est relationnelle.

\[
\boxed{
\Rac=\Rac(\Cadre,\Con,\text{conditions de sélection}).
}
\]

\subsection{3. La clôture locale n'est pas une clôture absolue}

Une racine peut être parfaitement stable dans un cadre sans constituer une
clôture du réel.

\[
\boxed{
\text{stabilité locale}\neq\text{absoluité ontologique}.
}
\]

Cette proposition permet de conserver la réalité des limites locales sans
les transformer en bornes métaphysiques ultimes.

\subsection{4. Pas d'extérieur absolu}

L'inracine ne doit pas être interprétée comme « l'extérieur de tout ». Si elle
était un extérieur absolu, elle deviendrait elle-même une totalité déterminée
par contraste avec un intérieur.

La conséquence correcte est plus faible et plus rigoureuse :
\[
\boxed{
\text{l'inracine est un dehors-relatif au cadre, non un dehors-absolu.}
}
\]

\subsection{5. Asymétrie fondamentale}

La structure ontologique n'est pas :
\[
\text{racine}\leftrightarrow\text{inracine}
\]
comme si les deux étaient deux catégories symétriques.

Elle est plutôt :
\[
\boxed{
\text{racine} = \text{détermination locale},
\qquad
\text{inracine} = \text{non-clôture relativement à cette détermination}.
}
\]

\subsection{6. Non-auto-suffisance de la racine}

Si une racine est une solution sous contraintes, elle n'est pas théoriquement
auto-suffisante :
\[
\Rac\in\Sol(\Cadre).
\]

Si les contraintes sont elles-mêmes des déterminations produites à un niveau
antérieur, on obtient :
\[
\mathfrak I_0
\to
\Con_0
\to
\Rac_0
\to
\Con_1
\to
\Rac_1
\to\cdots.
\]
Cette régression est un programme de recherche, non une preuve qu'elle est
nécessairement infinie dans tout modèle possible.

\section{Conséquences épistémologiques et perspectives}

\subsection{1. La perception comme cadrage}

Toute perception, mesure ou représentation sélectionne une classe de
caractéristiques accessibles depuis un cadre donné.

Ainsi :
\[
\boxed{
\text{observer}=\text{déterminer sous des contraintes d'accès et de représentation}.
}
\]

Cela explique pourquoi une description très précise peut rester locale.
La précision de la racine n'annule pas son inracine.

\subsection{2. Les perspectives sont relatives, pas arbitraires}

Deux cadres \(\Cadre_1\) et \(\Cadre_2\) peuvent produire deux racines
\(\Rac_1\) et \(\Rac_2\) différentes sans que l'une soit nécessairement une
« illusion ». Elles peuvent correspondre à des projections ou déterminations
différentes d'un même contexte ontologique.

Il faut toutefois conserver une condition de cohérence : deux perspectives ne
sont pas déclarées compatibles simplement parce qu'elles sont différentes.
Il faut examiner leurs domaines, contraintes et relations de traduction.

\subsection{3. Principe anti-réification}

\begin{proposition}[Principe anti-réification]
Le passage :
\[
\text{frontière d'un cadre}
\Rightarrow
\text{frontière de l'être}
\]
n'est pas logiquement valide sans prémisse supplémentaire.
\end{proposition}

Ce principe vaut pour le langage ordinaire, la perception, les modèles
scientifiques, les systèmes formels et les représentations informatiques.

\subsection{4. Connaître n'est pas totaliser}

Une connaissance peut être complète relativement à un cadre tout en restant
locale par rapport à un cadre plus englobant.

\[
\boxed{
\text{complétude relative}\neq\text{totalisation ontologique}.
}
\]

\section{Conséquences mathématiques}

\subsection{1. Deux niveaux à ne pas confondre}

La théorie respecte les mathématiques usuelles lorsqu'elles définissent
formellement la finitude, l'infinitude, les cardinalités ou les limites.
Elle ajoute une méta-question : quel statut ontologique attribue-t-on à ces
notions lorsqu'elles sont appliquées au réel ?

Ainsi :
\[
\boxed{
\text{fini dans une structure formelle}
\neq
\text{fini comme clôture ontologique ultime}.
}
\]

\subsection{2. Cadres comme structures de contraintes}

Un modèle mathématique \(M\) accompagné de contraintes \(C\) définit un espace
de solutions :
\[
\Sol_M(C)=\{x\in M\mid x\sat C\}.
\]
Une solution locale \(x\) peut être considérée comme une racine dans le
cadre \((M,C)\).

\subsection{3. Raffinement par contraintes}

Si \(C'\) renforce \(C\), c'est-à-dire si toute solution de \(C'\) satisfait
\(C\), alors :
\[
\boxed{
\Sol_M(C')\subseteq\Sol_M(C).
}
\]
La réduction de l'espace des solutions est donc une conséquence formelle
simple d'un renforcement des contraintes.

\subsection{4. Extension de cadre et non simple ajout de contraintes}

Une difficulté supplémentaire apparaît lorsqu'un nouveau niveau introduit
non seulement des contraintes supplémentaires, mais un langage, des variables
ou des structures inexistantes dans le cadre précédent.

Dans ce cas, il n'est pas toujours correct d'écrire :
\[
\Sol(C_{n+1})\subseteq\Sol(C_n).
\]
Il peut être nécessaire d'étendre le modèle :
\[
M_n\longrightarrow M_{n+1},
\]
puis de définir une relation de projection ou d'interprétation :
\[
\pi_{n+1,n}:M_{n+1}\to M_n.
\]

Cette distinction est importante pour ne pas transformer artificiellement
l'inracine en simple sous-ensemble caché d'un espace déjà fixé.

\subsection{5. Localité mathématique}

Une structure peut être finie dans une présentation et appartenir à une
structure plus riche dans une autre présentation. Il faut donc distinguer :
\begin{itemize}
    \item la finitude formelle interne à un système ;
    \item l'ouverture d'un modèle vers des extensions ;
    \item l'interprétation ontologique de ces structures.
\end{itemize}

\subsection{6. Le concept de racine en mathématiques}

Le concept de racine ne remplace pas toutes les notions mathématiques
existantes. Il fournit plutôt une couche méta-formelle :
\[
\boxed{
\text{racine} = \text{détermination admissible relativement à un cadre formel}.
}
\]

Cette couche peut s'appliquer à une solution d'équation, à un modèle, à un
état d'un système, à une structure satisfaisant des axiomes ou à un élément
sélectionné d'un espace de possibilités.

\section{Exemple : \(x+2=5\)}

On pose :
\[
C=\{x+2=5\}.
\]
Alors :
\[
\Sol(C)=\{3\}.
\]
Donc :
\[
\boxed{\Rac=3\in\Sol(C).}
\]

Le point ontologique n'est pas de nier que \(3\) est un nombre fini au sens
mathématique standard. Le point est plus précis : la propriété mathématique
de finitude de \(3\) ne suffit pas à établir que « fini » est une primitive
ontologique de la réalité.

\section{Exemple : le nombre 123}

On peut représenter :
\[
C_1\to\Rac_1=123,
\qquad
C_2\to\Rac_2=124.
\]

Le nombre 123 est une détermination dans un cadre. Il ne devient pas moins
réel parce qu'un autre cadre permet de distinguer 124.

La théorie introduit donc la différence entre :
\[
\text{« 123 est déterminé ici »}
\]
et :
\[
\text{« 123 est une borne ontologique absolue de tout ce qui peut être ».}
\]
La seconde proposition nécessite une prémisse supplémentaire qui n'est pas
contenue dans la première.

\section{Exemple : un univers décrit par 17 particules}

Supposons un modèle physique \(M\) avec la contrainte :
\[
C_{17}:\quad N=17.
\]
Alors :
\[
\Sol_M(C_{17})
\]
contient les états compatibles avec cette valeur.

Dans le cadre de la théorie, on peut écrire :
\[
\Rac_{17}\in\Sol_M(C_{17}).
\]
Cela n'implique pas :
\[
\text{« 17 est la limite ontologique ultime de tout réel possible ».}
\]

Si un cadre différent \(C_{18}\) est physiquement cohérent dans un modèle
étendu, on peut avoir :
\[
C_{18}:\quad N=18.
\]
Mais l'existence physique d'une telle possibilité n'est pas établie par la
seule métaphysique des racines et des inracines.

\begin{warning}[Possibilité conceptuelle vs possibilité physique]
Le fait qu'un nombre puisse être décrit ou qu'une contrainte puisse être
formellement modifiée n'établit pas que le monde physique réel autorise cette
modification.
\end{warning}

\section{Conséquences algorithmiques et computationnelles}

\subsection{1. Algorithme comme détermination locale}

Un algorithme peut être formalisé comme :
\[
\mathcal A:(\Cadre,E)\longrightarrow\Rac,
\]
où \(E\) désigne les données d'entrée et où la sortie \(\Rac\) satisfait des
contraintes spécifiées lorsque l'algorithme réussit.

On peut écrire :
\[
\Rac\in\Sol(\Cadre,E).
\]

\subsection{2. L'espace de recherche}

Une procédure algorithmique travaille sur un espace de candidats :
\[
\mathcal Q_0\supseteq\mathcal Q_1\supseteq\cdots\supseteq\mathcal Q_k.
\]
À chaque étape, les contraintes, tests ou règles de recherche peuvent réduire
la classe des candidats.

\begin{proposition}[Réduction de recherche]
Si chaque étape ajoute seulement des critères de rejet compatibles avec le
problème initial, alors :
\[
\mathcal Q_{n+1}\subseteq\mathcal Q_n.
\]
\end{proposition}

\subsection{3. L'inracine computationnelle}

Pour un algorithme donné, l'inracine peut être modélisée de manière
épistémique comme ce qui reste hors de la capacité de représentation, de
calcul ou de recherche du cadre courant.

On peut écrire :
\[
\Inrac^{\mathrm{comp}}(\mathcal A,\Cadre)
\]
pour désigner le résidu de non-clôture computationnelle.

Ce résidu peut provenir de plusieurs sources :
\[
\text{taille},\quad
\text{temps},\quad
\text{mémoire},\quad
\text{information manquante},\quad
\text{expressivité},\quad
\text{indécidabilité formelle}.
\]

\begin{warning}[Incalculable ne signifie pas inracine ontologique]
Une limite algorithmique particulière ne prouve pas qu'une limite ontologique
existe ou n'existe pas. L'algorithmique opère à l'intérieur d'un cadre donné.
\end{warning}

\subsection{4. Algorithme récursif de déplacement de cadre}

La théorie suggère un schéma général :
\[
\Cadre_n
\xrightarrow{\mathcal A_n}
\Rac_n
\xrightarrow{\text{diagnostic du résidu}}
\Inrac_n
\xrightarrow{\text{reformulation}}
\Cadre_{n+1}.
\]

Ce schéma ressemble structurellement aux méthodes d'extension ou de raffinement
de modèles : une solution est obtenue, puis l'analyse de ses limites motive un
nouveau cadre.

\begin{remark}
Ce processus n'implique aucune promesse de convergence vers une description
ultime. Au contraire, dans le cadre de l'axiome de non-totalisation, la
recherche d'une clôture exhaustive du fond n'est pas un objectif atteignable
par stipulation.
\end{remark}

\subsection{5. Existence, calculabilité, actualisation}

La théorie exige de distinguer au moins quatre notions :
\[
\begin{aligned}
&\text{admissible sous }\Cadre,\\
&\text{représentable dans }\Cadre,\\
&\text{calculable par }\mathcal A,\\
&\text{effectivement actualisé.}
\end{aligned}
\]

Aucune implication générale du type :
\[
\text{admissible}\Rightarrow\text{calculable}\Rightarrow\text{réalisé}
\]
n'est autorisée sans hypothèses supplémentaires.

\section{Conséquences physiques}

\subsection{1. Les états physiques comme racines}

Une formalisation physique peut être comprise comme :
\[
\Cadre_{\mathrm{phys}}
\to
\Sol(\Cadre_{\mathrm{phys}})
\to
\Rac_{\mathrm{phys}}.
\]
Les lois peuvent être interprétées comme structures de contraintes, mais cette
lecture reste une interprétation méta-physique tant qu'elle ne spécifie pas une
théorie dynamique particulière.

\subsection{2. Portée minimale}

La théorie n'établit pas, à elle seule :
\begin{itemize}
    \item que l'espace physique est infini ;
    \item que le temps est infini ;
    \item qu'il existe une infinité actuelle de particules ;
    \item que toute limite physique est illusoire ;
    \item qu'il existe une particule au-delà d'un nombre donné ;
    \item qu'une loi particulière constitue l'ultime mécanisme de contrainte.
\end{itemize}

\subsection{3. Ce que la théorie pourrait exiger d'une extension physique}

Pour devenir une théorie physique différenciable empiriquement, il faudrait
spécifier :
\begin{enumerate}
    \item des variables primitives ;
    \item des contraintes fondamentales ;
    \item une dynamique de génération, sélection ou actualisation ;
    \item une règle de correspondance avec les observations ;
    \item des prédictions nouvelles ou des restrictions observables.
\end{enumerate}

\section{Conséquences causales}

\subsection{Contrainte et cause}

Une contrainte ne doit pas être identifiée à une cause :
\[
\boxed{
\text{cause}\neq\text{contrainte}\neq\text{solution}.
}
\]

Une cause peut contribuer à produire une racine, alors qu'une contrainte peut
simplement définir ce qui est admissible.

\subsection{Causalité sous cadre}

Dans un modèle dynamique :
\[
\Rac_n(t)\xrightarrow{\text{dynamique}}\Rac_n(t+1).
\]
Dans un cadre plus large :
\[
\Rac_n\nextlev\Cadre_{n+1}.
\]
Il faut éviter de confondre les deux relations : la première peut être causale,
la seconde est une relation de changement de cadre.

\section{Conséquences structurelles et émergentes}

Une idée centrale de la Version 2 est que :
\[
\boxed{
\text{une racine à un niveau peut devenir une contrainte au niveau suivant}.
}
\]

Formellement :
\[
\Rac_n\nextlev\Con_{n+1}.
\]

Cette structure autorise une ontologie récursive : les déterminations déjà
obtenues ne sont pas seulement des résultats passifs ; elles peuvent devenir
des conditions structurantes pour de nouvelles déterminations.

\begin{consequence}[Émergence relationnelle]
L'émergence peut être représentée, dans ce cadre, comme l'apparition d'une
nouvelle racine dans un cadre dont les contraintes ne sont pas réductibles,
pour l'analyse considérée, au simple vocabulaire du niveau antérieur.
\end{consequence}

La théorie ne prétend pas que toute émergence est métaphysiquement irréductible.
Elle affirme seulement qu'un changement de cadre peut changer l'espace des
déterminations pertinentes.

\section{Conséquences modales : possible, représentable, réalisé}

Il faut distinguer :
\[
\boxed{
\begin{array}{lll}
\text{admissible} &:& D\in\Sol(\Cadre),\\
\text{représentable} &:& D\text{ peut être exprimé dans }\Cadre,\\
\text{possible} &:& D\text{ est compatible avec un cadre physiquement ou logiquement autorisé},\\
\text{réalisé} &:& D\text{ est effectivement actualisé dans le cadre considéré.}
\end{array}
}
\]

Ces notions peuvent coïncider dans certains modèles mais ne doivent pas être
fusionnées par défaut.

\subsection{Inracine et modalité}

L'inracine n'est pas l'ensemble de toutes les possibilités. La raison est
structurelle : définir « l'ensemble de toutes les possibilités » comme une
totalité exhaustive réintroduirait précisément le type de clôture que la notion
cherche à dépasser.

La formulation rigoureuse est donc :
\[
\boxed{
\text{l'inracine ouvre le cadre ; elle n'est pas la totalité de ce qui est possible.}
}
\]

\section{Conséquences linguistiques et sémantiques}

Les catégories utilisées dans le langage sont elles-mêmes des racines
conceptuelles.

On peut écrire :
\[
\text{contexte}
\to
\text{contraintes linguistiques}
\to
\text{concept}
\to
\text{racine sémantique}.
\]

Les mots « fini », « totalité », « extérieur », « objet » et même « réalité »
sont donc soumis à une précaution méthodologique : leur usage grammatical
ne garantit pas qu'ils désignent des bornes ontologiques absolues.

\begin{proposition}[Principe de non-réification sémantique]
La stabilité d'une distinction linguistique ne suffit pas à établir que la
structure distinguée constitue une séparation ontologique ultime.
\end{proposition}

\section{Conséquences informationnelles}

L'information elle-même peut être interprétée comme une détermination sous
contraintes.

Pour un agent \(A\), un état informationnel \(E_A\) est une représentation
relative à un canal, une mémoire, un langage et des ressources.

On peut écrire :
\[
E_A=E(A,\Cadre_A).
\]

L'inracine informationnelle correspond alors au résidu de ce qui n'est pas
capturé par cette représentation.

\begin{warning}[Résidu informationnel vs résidu ontologique]
Un résidu dans une base de données, une observation ou une mémoire est un
résidu informationnel. Il ne faut pas automatiquement l'identifier à un
résidu ontologique.
\end{warning}

\section{Conséquences sur la théorie des modèles et la connaissance scientifique}

Une théorie scientifique \(T_n\) fournit un cadre :
\[
T_n\to\Sol(T_n)\to\Rac_n.
\]

Lorsqu'une anomalie, une nouvelle échelle ou une nouvelle observation excède
le pouvoir explicatif de \(T_n\), plusieurs scénarios sont possibles :
\begin{enumerate}
    \item enrichir les contraintes ;
    \item modifier la structure du modèle ;
    \item étendre le vocabulaire ;
    \item changer de théorie.
\end{enumerate}

Dans le langage de la Version 2 : l'inracine d'un cadre scientifique peut
fonctionner comme un signal de non-clôture méthodologique. Cela ne signifie
pas que toute anomalie révèle une « nouvelle réalité » ; il faut un travail
empirique et théorique supplémentaire.

\section{Conséquences sur l'identité et la description des objets}

Supposons deux cadres \(\Cadre_1\) et \(\Cadre_2\) qui déterminent la même
entité sous des descriptions différentes. On peut avoir :
\[
\Rac_1=\mathrm{Desc}_1(O),
\qquad
\Rac_2=\mathrm{Desc}_2(O).
\]

La différence entre descriptions n'implique pas une multiplicité d'objets.
Mais il faut une relation d'interprétation \(\tau\) permettant de relier les
descriptions :
\[
\tau(\Rac_1)=\Rac_2
\]
ou, plus généralement, une correspondance partielle.

La théorie est donc compatible avec une ontologie où l'identité de l'objet
est distincte de chacune de ses racines descriptives.

\section{Conséquences sur la récursivité}

La structure récursive peut être distinguée en trois versions.

\subsection{1. Récursivité descriptive}
\[
\Cadre_n\to\Rac_n\to\Cadre_{n+1}.
\]
Le niveau supérieur est seulement un changement de description.

\subsection{2. Récursivité épistémique}
\[
\text{connaissance}_n
\to
\text{limite}_n
\to
\text{nouvelle connaissance}_{n+1}.
\]
L'agent augmente sa capacité de représentation.

\subsection{3. Récursivité ontologique}
\[
\Rac_n\to\Con_{n+1}\to\Rac_{n+1}.
\]
Une détermination réelle devient effectivement une condition de nouvelles
déterminations.

\begin{warning}[Ne pas confondre les trois récursivités]
Le fait qu'une structure puisse être décrite récursivement ne prouve pas que le
réel lui-même soit ontologiquement récursif. L'ontologie récursive est une
hypothèse supplémentaire.
\end{warning}

\section{Conséquences métaphysiques générales}

\subsection{1. Le fond n'est pas une solution finale}

Dans la théorie, la racine n'est pas l'aboutissement d'un calcul ayant trouvé
la « réponse totale » du réel.

\[
\boxed{
\Rac_n\neq\text{clôture ontologique finale}.
}
\]

\subsection{2. La connaissance est structurellement locale}

Toute représentation possède un cadre, donc une échelle, des critères de
pertinence et des exclusions.

\[
\boxed{
\text{toute représentation est locale relativement à des contraintes.}
}
\]

\subsection{3. L'inracine n'est pas un nouvel infini substantialisé}

Il serait contraire à la théorie de prendre :
\[
\Inrac=\text{un objet infiniment grand}
\]
ou :
\[
\Inrac=\text{un ensemble de tout ce qui échappe au cadre}.
\]

La notion doit rester relationnelle et ouverte.

\subsection{4. La distinction racine/inracine est elle-même relative}

Si une nouvelle racine est obtenue :
\[
\Inrac_n\to\Rac_{n+1},
\]
alors ce qui était inracine au niveau \(n\) devient partie d'une nouvelle
racine au niveau \(n+1\). Mais \(\Rac_{n+1}\) n'épuise pas nécessairement
l'inracine dont elle provient.

Ainsi :
\[
\boxed{
\text{ce qui devient racine à un niveau peut conserver une inracine à un niveau supérieur.}
}
\]

\section{Une formalisation plus fine de la notion de « global »}

L'expression « l'inracine est la globalité » doit être précisée pour éviter une
contradiction entre globalité et non-totalisation.

\begin{definition}[Globalité relative]
Une structure est dite \emph{globalement relative} à un cadre \(\Cadre\) si
elle excède les frontières locales de \(\Cadre\) tout en restant dépourvue
d'une prétention à épuiser toute réalité.
\end{definition}

On peut alors écrire :
\[
\boxed{
\Inrac(\Cadre)=\text{globalité relative du cadre, sans totalisation absolue}.
}
\]

Le mot « global » signifie donc « englobant pour ce cadre-ci », et non
« totalité absolue de tout ce qui existe ».

\section{Le problème de « l'inracine contraint la racine »}

La formulation initiale :
\[
\text{l'infini contraint le fini}
\]
était difficile parce qu'une contrainte ordinaire réduit un espace de solutions,
alors que l'infini était défini comme non-limitable.

La Version 2 remplace cette phrase par :
\[
\boxed{
\text{l'inracine conditionne les limites locales de la racine.}
}
\]

Une formulation encore plus stricte est :
\[
\boxed{
\text{la racine n'est locale que relativement à l'inracine de son cadre.}
}
\]

Cette formulation ne prétend plus que l'inracine exerce une contrainte
identique aux contraintes locales. Elle introduit une relation de fondement,
de contexte ou de non-clôture.

\section{Question fondamentale : qu'est-ce qui contraint l'inracine ?}

Dans la théorie, demander « qu'est-ce qui contraint l'inracine fondamentale ? »
peut conduire à une contradiction si l'inracine fondamentale est précisément
le principe non clôturable qui n'est pas lui-même dérivé d'une racine supérieure.

La réponse interne de la théorie est donc :
\[
\boxed{
\mathfrak I_\ast\text{ n'est pas posé comme solution d'une contrainte supérieure.}
}
\]

Cela ne constitue pas une démonstration indépendante ; c'est un engagement
axiomatique de la théorie.

\section{Statut de la définition humaine de l'infini}

Le terme « infini » reste historiquement et philosophiquement à l'origine de
l'intuition de la théorie, mais la Version 2 déplace l'accent vers
l'inracine.

Une représentation humaine :
\[
D_{\mathrm{hum}}(\text{infini})
\]
est une racine conceptuelle. Elle ne doit pas être identifiée au fond qu'elle
cherche à désigner.

On peut alors écrire :
\[
\boxed{
\text{toute conception de l'infini est une racine conceptuelle et porte son inracine.}
}
\]

La théorie évite ainsi deux erreurs opposées :
\begin{enumerate}
    \item prétendre que notre définition épuiserait le non-clôturable ;
    \item conclure que, parce qu'une conception est imparfaite, elle est nécessairement fausse.
\end{enumerate}

\section{Limites logiques de la formalisation}

La formalisation proposée ne doit pas être présentée comme une déduction
mathématique complète d'une métaphysique depuis zéro. Elle contient des axiomes
philosophiques.

En particulier, les propositions suivantes sont des engagements du cadre :
\[
\begin{aligned}
&\text{(A1)}&&\mathfrak I_\ast\text{ est fondamental},\\
&\text{(A2)}&&\mathfrak I_\ast\text{ n'est pas totalisable},\\
&\text{(A3)}&&\Rac\in\Sol(\Cadre),\\
&\text{(A4)}&&\text{les niveaux peuvent se relier récursivement}.
\end{aligned}
\]

À partir de ces axiomes, les résultats mathématiques locaux tels que
\(\Sol(C')\subseteq\Sol(C)\) sont démontrables dans les cadres formels
correspondants. En revanche, les affirmations physiques ou métaphysiques
universelles exigent des principes de pont supplémentaires.

\section{Questions ouvertes structurantes}

\begin{question}[Origine des contraintes]
Les contraintes sont-elles elles-mêmes fondamentales, émergentes ou produites
par des racines antérieures ?
\end{question}

\begin{question}[Sélection]
Si \(|\Sol(\Cadre)|>1\), qu'est-ce qui sélectionne une racine particulière ?
\end{question}

\begin{question}[Actualisation]
Quelle relation distingue une solution admissible d'une solution effectivement
réalisée ?
\end{question}

\begin{question}[Pont physique]
Quelles contraintes observables permettraient de transformer cette ontologie
en théorie physique ?
\end{question}

\begin{question}[Pont mathématique]
Existe-t-il une structure mathématique non triviale représentant la relation
racine/inracine sans réifier l'inracine en un ensemble total ?
\end{question}

\begin{question}[Convergence]
Une suite de raffinements :
\[
\Cadre_0\to\Cadre_1\to\Cadre_2\to\cdots
\]
peut-elle converger vers une limite locale sans jamais produire une clôture
ontologique absolue ? Sous quelles notions de convergence ?
\end{question}

\begin{question}[Invariance]
Quelles propriétés d'une racine restent invariantes lorsque le cadre change ?
Cette question est nécessaire si l'on veut distinguer variation de description
et variation ontologique réelle.
\end{question}

\section{Conséquences et critères de rigueur}

Pour qu'une future théorie fondée sur l'inracine reste fidèle à son noyau,
elle devrait satisfaire au minimum les critères suivants :
\begin{enumerate}
    \item ne pas transformer l'inracine en ensemble de « tout ce qu'on ne connaît pas » ;
    \item ne pas traiter une racine locale comme une frontière absolue ;
    \item distinguer contraintes, causes, représentations et réalisations ;
    \item distinguer extension de cadre et simple ajout de contraintes ;
    \item ne pas confondre possibilité conceptuelle et possibilité physique ;
    \item ne pas supposer qu'une récursivité descriptive prouve une récursivité ontologique ;
    \item maintenir séparés les niveaux ontologique, épistémique, mathématique et computationnel ;
    \item considérer l'absence de totalisation comme un axiome métaphysique, et non comme un théorème mathématique automatique.
\end{enumerate}

\section{Tableau général des conséquences}

\renewcommand{\arraystretch}{1.18}
\begin{longtable}{>{\raggedright\arraybackslash}p{0.18\textwidth}
                  >{\raggedright\arraybackslash}p{0.34\textwidth}
                  >{\raggedright\arraybackslash}p{0.38\textwidth}}
\toprule
\textbf{Niveau} & \textbf{Réinterprétation} & \textbf{Conséquence}\\
\midrule
\endfirsthead
\toprule
\textbf{Niveau} & \textbf{Réinterprétation} & \textbf{Conséquence}\\
\midrule
\endhead
Ontologique & Le fond est non clôturable ; la racine est locale. & Une borne locale n'est pas automatiquement une clôture absolue.\\
Épistémologique & Connaître signifie déterminer dans un cadre. & Toute connaissance est située relativement à des contraintes d'accès et de représentation.\\
Perspective & Une perspective est un cadre local. & Deux perspectives peuvent être différentes sans que la différence suffise à établir une contradiction ou une identité.\\
Mathématique & Une solution est une détermination satisfaisant des contraintes formelles. & Les espaces de solutions se réduisent sous ajout de contraintes ; les extensions de modèle demandent des relations de traduction.\\
Algorithmique & Calculer revient à construire ou sélectionner une racine dans un espace de recherche. & La calculabilité est toujours relative à un cadre, des ressources et une procédure.\\
Physique & Un état physique peut être représenté comme solution de contraintes. & Une ontologie des racines n'est pas encore une théorie physique.\\
Causal & Les causes et les contraintes sont distinctes. & L'admissibilité n'est pas automatiquement la causalité.\\
Modal & Possible, représentable et réalisé sont distincts. & L'existence d'une solution formelle ne garantit pas son actualisation.\\
Linguistique & Les concepts sont des déterminations. & Une distinction grammaticale ne prouve pas une frontière ontologique.\\
Informationnel & Toute représentation possède des limites de canal, mémoire et langage. & Un résidu informationnel n'est pas nécessairement un résidu ontologique.\\
Structurel & Une racine peut devenir contrainte. & Une ontologie récursive des déterminations est possible comme extension de la théorie.\\
Scientifique & Un modèle est une racine théorique sous contraintes. & L'inracine peut signaler une non-clôture méthodologique, sans constituer à elle seule une preuve empirique.\\
Métaphysique & L'inracine n'est ni objet maximal ni totalité. & La théorie évite de réintroduire la totalité par le vocabulaire du « dehors ».\\
\bottomrule
\end{longtable}

\section{Synthèse formelle minimale}

Le système complet peut être résumé par les relations suivantes.

\subsection*{Fondement}
\[
\boxed{
\mathfrak I_\ast
\text{ est non clôturable et n'est pas une racine d'un cadre supérieur.}
}
\]

\subsection*{Localisation}
\[
\boxed{
\Cadre_n
\to
\Sol(\Cadre_n)
\to
\Rac_n.
}
\]

\subsection*{Inracine}
\[
\boxed{
\Rac_n,\Cadre_n
\to
\Inrac_n
=
\text{dépassement non clôturé de }\Cadre_n.
}
\]

\subsection*{Récursivité}
\[
\boxed{
\Inrac_n
\nextlev
\Cadre_{n+1}
\to
\Sol(\Cadre_{n+1})
\to
\Rac_{n+1}
\to
\Inrac_{n+1}
\to\cdots
}
\]

\subsection*{Fini relatif}
\[
\boxed{
\mathrm{Fini}_{\Cadre_n}(\Rac_n)
\not\Rightarrow
\text{clôture ontologique absolue}.
}
\]

\subsection*{Principe central}
\[
\boxed{
\textbf{Toute racine porte son inracine.}
}
\]

\section{Conclusion}

La Version 2 ne dit pas simplement que « tout est infini ». Elle propose une
architecture conceptuelle plus précise :

\[
\boxed{\begin{gathered}
\text{un fond non clôturable}
\to
\text{des cadres}
\to
\text{des contraintes}
\to
\text{des solutions}\\
\to
\text{des racines locales}
\to
\text{des inracines}
\to
\text{de nouveaux cadres}
\to\cdots
\end{gathered}}
\]

Le concept d'\emph{inracine} permet de reformuler l'intuition initiale sans
faire de l'infini une « force qui limite » au sens ordinaire. L'inracine est
plutôt le nom donné au fait structurel qu'une racine n'est locale que
relativement à un cadre, et qu'un tel cadre ne constitue pas, par lui-même,
une clôture ontologique ultime.

Cette reformulation conserve les idées fondamentales développées auparavant :
le fond non-limitable est premier ; les déterminations apparaissent relativement
à des contraintes ; ce que nous appelons « fini » peut être compris comme une
propriété secondaire de solutions locales ; une solution n'est pas
nécessairement une réalisation ; et une détermination obtenue à un niveau peut
servir de contrainte ou de point d'appui à un niveau ultérieur.

Le déplacement essentiel est maintenant le suivant :
\[
\boxed{
\text{infini vs fini}
\quad\longrightarrow\quad
\text{racine vs inracine}.
}
\]

Le mot « inracine » ne supprime donc pas le problème de l'infini ; il en change
le point d'accès théorique. Au lieu de chercher à représenter directement une
totalité non totalisable, la théorie décrit la relation entre une détermination
locale et ce qui excède sa clôture. L'inracine n'est ainsi ni un objet caché,
ni l'ensemble de toutes les possibilités, ni un simple synonyme de notre
ignorance. C'est une relation de non-clôture.

La théorie demeure néanmoins une métaphysique axiomatique. Ses conséquences
les plus fortes --- existence physique du fond non clôturable, caractère réel
d'une récursivité ontologique, nature des contraintes fondamentales,
mécanisme de sélection des solutions --- demandent des arguments ou des
axiomes supplémentaires. Sa rigueur exige donc de distinguer en permanence ce
qui est \emph{défini}, ce qui est \emph{déduit}, ce qui est \emph{proposé} et ce
qui reste \emph{ouvert}.

\end{document}
